\section{Candor}

The central claim is a conjecture. It is not yet a completed result. Everything in this paper is an in-silico computation on the
substrate rather than a measurement against nature. The gaps are stated plainly here, one short paragraph each, so they are
not buried under the firm results.

\textbf{A field theory is not yet built.} The exact result is that the discrete divergence of the conserved charge is zero.
A full conserved-current field theory needs a named continuum field, an action, a stress tensor, and equations of motion.
The map from the lattice sum rule to a specific field with a stress tensor that sources curvature is unbuilt. Continuity is
shown, but the field theory it implies has not been built.

\textbf{Exact Lorentz invariance on a fixed lattice.} This is the hardest open problem. A fixed deterministic lattice is
standardly expected to pick a preferred frame, the frame in which it sits still. Causal sets reach exact Lorentz invariance
only by random sprinkling, the move vibe forbids \cite{bombelli1987, dowker2004}. Vibe's reply is a doubly-special
relativity with two invariants, the speed of light and the lattice spacing, and a discrete $F_4$ that restores to a
continuous symmetry in the infrared. Shown so far: a ballistic light cone, a conserved momentum current with sharp speed
one, and a doubly-special dispersion that recovers Einstein at low energy. Still open: full exact Lorentz from the fixed lattice.
If the lattice left a permanent preferred frame at low order, precision astrophysics would already have seen the Lorentz
violation, and it has not \cite{thooft2016}.

\textbf{Three chiral generations.} Two forced threefolds exist on the substrate, but their identification with the three
observed generations fails the naive match. The Distler-Garibaldi no-go forbids fitting three chiral generations into $E_8$
in the continuum without mirror fermions that nature lacks, and it constrains every algebra-spine program alike \cite{distler2010,
boyle2020}. Vibe's only structural difference is discreteness, which may evade the continuum no-go but has not been shown
to. The naive triality reading gives a vector plus two chiralities, a negative result. The escape is open.

\textbf{The single measurement outcome.} Vibe reproduces the Born rule and the Tsirelson bound from a deterministic base
with no injected randomness, and every emergent mixed state is the marginal of a pure reversible whole. What it does not
deliver is why a single definite outcome is recorded on one run. The measurement step is reached as statistics, not as the
selection of one result. The single-outcome problem is open here, as it is everywhere.

\textbf{Selves as causally real in the lossy middle.} The closed reversible base favors causal reduction, so selves must
come from lossy coarse-graining or open subsystems. Whether vibe's selves are irreducible high-order causes, in the sense
of causal emergence \cite{hoel2013}, is unmeasured. The bound self-repairing body still needs one discrete ingredient that
has not been isolated. Open.

\textbf{Emergent locality at scale.} A coarse-grained theory is only physical if its effective interactions stay local once
averaged \cite{konopka2008}. That the substrate keeps locality at scale, rather than developing the disordered nonlocality
a graph theory can suffer, is a well-posed test that has not been run to completion. Open.

\textbf{The other side of the bridge is itself unconfirmed.} The continuum program vibe is measured against is unreviewed,
unreplicated, and single-author \cite{herbert2025chronoflux}. The bridge takes its principles, targets, and methods, and
leaves the rest, but a parallel to an unconfirmed program is a weaker anchor than a parallel to an established one. This is
stated plainly.

A page claiming every problem is solved would be worth less than this list of open ones. Each open problem narrows the search
and points to what to test next. The thesis is a conjecture that the firm results support but do not prove.
